\documentclass[11pt,a4paper]{article}

% Template visual Matriz Aprovação. Compile com XeLaTeX ou LuaLaTeX.
\usepackage{fontspec}
\usepackage{polyglossia}
\setmainlanguage{portuguese}
\IfFontExistsTF{Aptos}{
  \setmainfont{Aptos}
  \setsansfont{Aptos}
}{
  \IfFontExistsTF{Calibri}{
    \setmainfont{Calibri}
    \setsansfont{Calibri}
  }{
    \IfFontExistsTF{TeX Gyre Heros}{
      \setmainfont{TeX Gyre Heros}
      \setsansfont{TeX Gyre Heros}
    }{
      \setmainfont{Latin Modern Sans}
      \setsansfont{Latin Modern Sans}
    }
  }
}
\usepackage[a4paper,margin=2cm,headheight=16pt]{geometry}
\usepackage{graphicx}
\usepackage{xcolor}
\usepackage{tikz}
\usepackage{tcolorbox}
\usepackage{enumitem}
\usepackage{booktabs}
\usepackage{tabularx}
\usepackage{amssymb}
\usepackage{fancyhdr}
\usepackage{titlesec}
\usepackage{hyperref}

\definecolor{matrizlime}{HTML}{CBFF4D}
\definecolor{matrizink}{HTML}{101313}
\definecolor{matrizmuted}{HTML}{596363}
\definecolor{matrizline}{HTML}{DDE3DF}
\definecolor{matrizsoft}{HTML}{F3F7EC}

\hypersetup{colorlinks=true,linkcolor=matrizink,urlcolor=matrizink}
\pagestyle{fancy}
\fancyhf{}
\lhead{\textsf{\small MATRIZ APROVAÇÃO}}
\rhead{\textsf{\small APOSTILA}}
\cfoot{\textsf{\small \thepage}}
\renewcommand{\headrulewidth}{0.4pt}
\renewcommand{\headrule}{\hbox to\headwidth{\color{matrizline}\leaders\hrule height \headrulewidth\hfill}}

\titleformat{\section}{\Large\bfseries\color{matrizink}}{\thesection}{0.6em}{}
\titleformat{\subsection}{\large\bfseries\color{matrizink}}{\thesubsection}{0.6em}{}
\setlist{nosep,leftmargin=*}

\newtcolorbox{foco}{colback=matrizsoft,colframe=matrizlime,boxrule=1pt,arc=3pt,left=10pt,right=10pt,top=8pt,bottom=8pt}
\newtcolorbox{lei}{colback=matrizink,colframe=matrizink,coltext=white,boxrule=0pt,arc=3pt,left=10pt,right=10pt,top=8pt,bottom=8pt}
\newtcolorbox{alerta}{colback=white,colframe=matrizline,boxrule=0.6pt,arc=3pt,left=10pt,right=10pt,top=8pt,bottom=8pt}

% Logo usada na capa. Caminho relativo à pasta onde o .tex é compilado
% (materiais/<disciplina>/). Ajuste se a estrutura mudar.
\newcommand{\logomatriz}{../../public/logo.png}

\begin{document}

\begin{titlepage}
  \thispagestyle{empty}
  \begin{center}
    \includegraphics[width=0.55\linewidth]{\logomatriz}\par
  \end{center}
  \vspace{2.2cm}
  {\sffamily\fontsize{28}{32}\selectfont\bfseries Apostila — Português: Interpretação de texto\par}
  \vspace{0.4cm}
  {\sffamily\large concursos gerais, OAB, carreiras militares e ENEM\par}
  \vfill
  \begin{foco}
    \textsf{\textbf{Foco da apostila}}\\[3pt]
    a alternativa correta precisa ser comprovada pelo texto. Uma alternativa apenas possível, provável ou compatível com o conhecimento de mundo não basta.
  \end{foco}
  \vspace{0.7cm}
  {\sffamily\small Atualizado em 16 de agosto de 2026\quad ·\quad adaptação necessária ao edital e à banca\par}
  \vspace{1cm}
  {\sffamily\small matrizaprova.com\par}
\end{titlepage}

\tableofcontents
\newpage

% Conteúdo gerado entra aqui.
\section*{O que cai}

Interpretação não é procurar uma frase que pareça correta. É reconstruir o sentido produzido pelo texto e verificar se a alternativa mantém esse sentido. As cobranças mais frequentes são:

\begin{enumerate}
  \item identificação da ideia central e da finalidade do texto;
  \item distinção entre informação explícita e inferência autorizada;
  \item relação entre partes do texto, especialmente causa, oposição, conclusão e
\end{enumerate}

   explicação;

\begin{enumerate}
  \item efeito de sentido de palavras, expressões, pontuação e conectivos;
  \item compatibilidade entre uma afirmação da alternativa e o texto inteiro.
\end{enumerate}

\textbf{Regra de prova:} a alternativa correta precisa ser comprovada pelo texto. Uma alternativa apenas possível, provável ou compatível com o conhecimento de mundo não basta.

\section*{Teoria essencial}

\subsection*{1. Tema, ideia central e finalidade}

O \textbf{tema} é o assunto geral. A \textbf{ideia central} é o que o texto afirma sobre esse assunto. A \textbf{finalidade} é o objetivo comunicativo: informar, criticar, convencer, explicar, narrar, orientar, ironizar ou provocar reflexão.

Exemplo:

\begin{alerta}
O uso de aplicativos facilitou o acesso a serviços públicos, mas também transferiu ao cidadão a responsabilidade de resolver falhas que antes eram atendidas presencialmente.
\end{alerta}

\begin{itemize}
  \item Tema: uso de aplicativos em serviços públicos.
  \item Ideia central: a digitalização facilita o acesso, mas pode transferir custos e responsabilidades ao cidadão.
  \item Finalidade provável: analisar criticamente uma consequência da digitalização.
\end{itemize}

Para localizar a ideia central, observe título, início, conclusão e repetições. Uma informação repetida com palavras diferentes costuma indicar o eixo do texto.

\subsection*{2. Informação explícita e inferência}

Informação \textbf{explícita} está declarada. Inferência é uma conclusão necessária obtida pela combinação de informações do próprio texto.

Se o texto afirma que "a procura cresceu após a redução do preço", pode-se inferir que existe relação temporal entre redução e crescimento. Não se pode, sem outra informação, afirmar que a redução foi a única causa, que o produto ficou melhor ou que o crescimento continuará.

Use o teste da prova textual:

\begin{enumerate}
  \item localize a passagem relacionada à alternativa;
  \item substitua a alternativa por uma pergunta: o texto permite essa conclusão?
  \item elimine palavras absolutas que extrapolem o que foi escrito.
\end{enumerate}

\subsection*{3. Pressuposto e subentendido}

O \textbf{pressuposto} é uma informação acionada por uma palavra ou construção. Em "a empresa voltou a contratar", o verbo \textit{voltar} pressupõe que ela já contratou antes. O pressuposto permanece mesmo quando a frase é negada: "a empresa não voltou a contratar" ainda pressupõe contratação anterior.

O \textbf{subentendido} depende mais do contexto e da intenção. Em uma situação em que alguém pergunta "Você vai mesmo entregar isso hoje?", o advérbio \textit{mesmo} pode sugerir dúvida, embora a dúvida não seja afirmada diretamente.

\subsection*{4. Conectivos e relações de sentido}

Conectivos organizam a argumentação. Não basta decorar a palavra: é preciso observar a relação entre as orações.

\begin{tabularx}{\linewidth}{llX}
\toprule
 \textbf{Relação} & \textbf{Exemplos} & \textbf{Efeito} \\
\midrule
 adição & e, também, além disso & acrescenta informação \\
 oposição & mas, porém, contudo & contrapõe ideias \\
 causa & porque, visto que, já que & apresenta motivo \\
 consequência & portanto, por isso, assim & apresenta resultado \\
 condição & se, caso, desde que & estabelece hipótese \\
 concessão & embora, ainda que & admite obstáculo sem impedir o fato \\
 conclusão & logo, assim, desse modo & fecha o raciocínio \\
\bottomrule
\end{tabularx}


Trocar um conectivo pode alterar completamente a relação lógica. Em questões de reescrita, compare não apenas o vocabulário, mas também causa, tempo, condição, oposição e conclusão.

\subsection*{5. Referência e coesão}

Pronomes e expressões retomam ou antecipam informações.

\begin{alerta}
A diretora conversou com a equipe antes de \textbf{ela} divulgar o resultado.
\end{alerta}

O pronome pode gerar ambiguidade se houver mais de um antecedente possível. Para resolver a questão, procure concordância, posição no período e sentido global. A referência não deve ser escolhida apenas pela palavra mais próxima.

Também há elipse: um termo fica oculto, mas é recuperado pelo contexto. Em "João preferiu o relatório; Maria, o parecer", o verbo \textit{preferiu} está oculto na segunda oração.

\subsection*{6. Vocabulário no contexto}

O sentido de uma palavra depende da situação. \textit{Frio} pode indicar temperatura, distanciamento emocional ou cálculo racional. Em prova, substitua a palavra pela opção indicada e releia o período inteiro. A substituição só é válida se preservar sentido, classe gramatical e tom.

Ironia merece atenção especial: o autor pode afirmar literalmente uma coisa para sugerir a crítica do contrário. A identificação depende de contraste, exagero, contexto e sinais de avaliação.

\subsection*{7. Argumentação e ponto de vista}

Textos argumentativos apresentam uma tese e razões para sustentá-la. Pergunte:

\begin{itemize}
  \item qual posição o autor defende?
  \item quais fatos ou exemplos funcionam como evidência?
  \item há oposição a uma opinião anterior?
  \item a conclusão é explícita ou decorre do conjunto?
\end{itemize}

Não confunda opinião do autor com opinião citada. Expressões como "segundo", "para os críticos" e "alguns defendem" introduzem vozes de terceiros.

\section*{Como a banca cobra}

\begin{itemize}
  \item \textbf{Certo/errado:} uma palavra absoluta como \textit{sempre}, \textit{nunca}, \textit{somente} ou \textit{necessariamente} pode tornar a afirmação incorreta.
  \item \textbf{Reescrita:} confira se a troca preserva sentido, relação lógica, tempo e referência dos pronomes.
  \item \textbf{Inferência:} escolha apenas a conclusão obrigatória pelo texto, não a que parece razoável fora dele.
  \item \textbf{Título:} o título pode resumir o tema, mas não substitui a ideia central.
  \item \textbf{Vocabulário:} o sentido contextual prevalece sobre o sentido de dicionário.
\end{itemize}

\subsection*{Exemplos resolvidos}

\subsubsection*{Exemplo 1}

Texto: "A biblioteca ampliou o horário, mas o número de empréstimos permaneceu estável. A direção decidiu, então, investir na divulgação do acervo."

\textbf{Pergunta:} a decisão da direção decorreu da percepção de que:

A) o horário ampliado reduziu o interesse dos leitores. \\ B) o acervo era insuficiente. \\ C) a ampliação do horário não aumentou os empréstimos. \\ D) os leitores preferiam comprar livros. \\ E) a biblioteca deveria reduzir o atendimento.

\begin{alerta}
\textbf{Gabarito: C.} O texto liga a estabilidade dos empréstimos à decisão de investir em divulgação. As demais alternativas acrescentam informações não apresentadas.
\end{alerta}

\subsubsection*{Exemplo 2}

Texto: "Embora a ferramenta reduza o tempo de busca, seu uso exige critérios: resultados rápidos não são necessariamente resultados confiáveis."

O conectivo \textit{embora} introduz uma ideia de:

A) causa. B) condição. C) conclusão. D) concessão. E) finalidade.

\begin{alerta}
\textbf{Gabarito: D.} A oração admite uma vantagem da ferramenta, mas apresenta uma restrição que impede a conclusão de que rapidez garante confiabilidade.
\end{alerta}

\section*{Quadro-resumo}

\begin{tabularx}{\linewidth}{lX}
\toprule
 \textbf{Pergunta} & \textbf{Procedimento} \\
\midrule
 Qual é o tema? & assunto geral, sem a opinião completa \\
 Qual é a ideia central? & afirmação principal sobre o tema \\
 O que é inferência? & conclusão necessária a partir do texto \\
 Como avaliar conectivo? & observe a relação entre as orações \\
 Como avaliar vocabulário? & substitua e releia o período inteiro \\
 Como detectar extrapolação? & procure informação que não está comprovada \\
 Como lidar com "sempre"? & verifique se o texto realmente autoriza totalidade \\
\bottomrule
\end{tabularx}


\section*{Questões de fixação}

\subsection*{Questão 1 — (questão inédita)}

"A automatização não elimina a necessidade de trabalho humano; modifica as habilidades exigidas e desloca o esforço para tarefas de análise."

A ideia central é que a automatização:

A) elimina trabalhos de análise. \\ B) torna desnecessárias as habilidades humanas. \\ C) transforma as habilidades necessárias ao trabalho. \\ D) reduz todo esforço profissional. \\ E) impede a análise de dados.

\begin{alerta}
\textbf{Gabarito: C.} A alternativa reproduz a relação expressa por "modifica as habilidades exigidas". As demais apresentam generalizações que contradizem ou extrapolam o texto.
\end{alerta}

\subsection*{Questão 2 — (questão inédita)}

Em "O projeto foi retomado", o verbo \textit{retomar} pressupõe que o projeto:

A) foi concluído. B) havia sido iniciado antes. C) foi rejeitado. D) será alterado. E) é obrigatório.

\begin{alerta}
\textbf{Gabarito: B.} Retomar significa voltar a realizar algo interrompido ou já iniciado.
\end{alerta}

\subsection*{Questão 3 — (questão inédita)}

No período "O relatório é extenso; portanto, sua leitura exige planejamento", o termo \textit{portanto} estabelece relação de:

A) oposição. B) causa. C) condição. D) conclusão. E) concessão.

\begin{alerta}
\textbf{Gabarito: D.} A segunda oração apresenta a conclusão decorrente da primeira.
\end{alerta}

\subsection*{Questão 4 — (questão inédita)}

Se um texto afirma que uma medida "pode reduzir custos", é correto concluir que a medida:

A) sempre reduz custos. B) já reduziu custos. C) necessariamente será adotada. D) tem potencial de reduzir custos. E) reduz custos sem efeitos colaterais.

\begin{alerta}
\textbf{Gabarito: D.} O verbo modal \textit{pode} indica possibilidade, não certeza.
\end{alerta}

\subsection*{Questão 5 — (questão inédita)}

Em uma questão de reescrita, a troca de "porém" por "portanto" é inadequada quando:

A) há uma conclusão explícita. B) há adição de argumentos. C) o período apresenta oposição. D) o período apresenta causa. E) o texto termina uma explicação.

\begin{alerta}
\textbf{Gabarito: C.} \textit{Porém} introduz oposição; \textit{portanto} introduz conclusão. A troca altera a relação lógica original.
\end{alerta}

\section*{Revisão final}

\begin{itemize}
  \item[$\square$] Separei tema de ideia central.
  \item[$\square$] Procurei a prova textual antes de aceitar uma inferência.
  \item[$\square$] Identifiquei a relação dos conectivos.
  \item[$\square$] Verifiquei pressupostos e referências pronominais.
  \item[$\square$] Analisei palavras no contexto, não isoladamente.
  \item[$\square$] Desconfiei de alternativas absolutas ou com informação nova.
  \item[$\square$] Conferi se a reescrita preserva sentido e lógica.
  \item[$\square$] Resolvi as questões sem usar conhecimento externo ao texto.
\end{itemize}

\end{document}
